\section{Lezione 18} \hfill 14/04/2026
\subsection{Vivai forestali}
Il materiale forestale da utilizzare dev'essere lavorato: per prima cosa va trovato il seme da seminare in vivaio.\\
Non va bene qualunque seme da piantare: le specie arboree generano razze diverse a seconda del luogo in cui crescono (es. ai bordi). Il seme deve provenire da boschi molto simili a quelli che si ha intenzione di rinnovare.\\
Il commercio dei semi non è libero, ma vincolato. Se il seme non viene acquistato, si può fare ciò che si vuole. Nel caso si andasse a commercializzarlo (comprare o vendere), occorre seguire delle leggi. In Italia, occorre raccogliere il seme nei ``boschi da seme'': fino a poco tempo fa la raccolta e commercio iniziava da parte dello Stato.\\
Il registro nazionale dei boschi da seme è un concetto vetusto, relativo al vecchio commercio di semi da conifere: esisteva solamente un bosco per la raccolta di semi da faggio (che era il Cansiglio).\\
Negli anni '90, è iniziata l'arboricoltura da legno, alternativa alla pioppicoltura (es. ciliegio, rovere, aceri): per queste specie non esistevano boschi da seme. Per ovviare a questo problema, le Regioni hanno codificato i boschi o piante da seme per l'arboricoltura (e di conseguenza anche per la selvicoltura).\\
%Occorre valutare bene la provenienza del seme da impianto, poiché dove provenire da un bosco simile a quello che si ha intenzione di rimpiantare.\\
Il seme dev'essere raccolto, e questo genera dei problemi per il materiale da utilizzare: ci sono semi facili da raccogliere (es. ghiande, ciliegie) e per altri è più difficile (es. abete bianco, il quale cono si sfalda sull'albero).\\
Il bosco ha solitamente degli alberi con fenotipi (e forse anche genotipi) migliori: chi raccoglie i semi dovrebbe preferire gli individui migliori (in verità tende a scegliere quelli più comodi). In questo modo c'è una riduzione della biodiversità: si raccolgono i semi da poche piante, non essendo sicuri che siano della migliore qualità (cosa che naturalmente non avviene).\\
La raccolta del seme, solitamente, segue il ciclo delle annate di pasciona (es. il castagno lo ha tutti gli anni, le querce ogni 2-3, il faggio ogni 5-6, il larice ogni 10-12).\\
I semi, dal punto di vista della conservazione sono considerati:
\begin{itemize}[nosep]
    \item recalcitranti: non si conservano facilmente perché hanno molta acqua (es. ghiande). Questi vengono piantati quasi subito, solitamente nella stagione successiva (nel frattempo vengono conservati nella sabbia);
    \item ortodossi.
\end{itemize}
Il modo migliore per conservare il seme oggi è il congelamento. Certi semi prima di essere piantati, devono essere congelati per farli attivare (poiché possiedono la strategia della vernalizzazione).\\
I semi vengono poi puliti: si eliminano foglie, rametti, pietruzze. Certe volte, prima di essere piantato viene trattato: l'assenza delle naturali condizioni ambientali rendono il tegumento duro. Perciò, vengono usati sistemi chimici o meccanici per ammollire il tegumento (es. acidi). Sempre più spesso i semi vengono inconfettati, ovvero ricoperti con sostanze speciali (es. insetticidi, anticrittogamici).\\
Infine, i semi vengono tolti dai luoghi di conservazione e seminati in vivaio.\\
È intrinseco che per ogni specie, una certa quota di semi non è in grado di germinare: normalmente, per una specie normale, il 6-8\% dei semi non germina (es. abete rosso). In altre specie (es. larice), la quota di ingerminabilità può raggiungere il 50\%.\\
Ci sono due tipologie di vivaio:
\begin{itemize}[nosep]
    \item per materiale a radici nude: si lavora in campo, nel quale vengono ottenute delle aiuole di coltura (strette e lunghe). Nella aiuola il seme viene seminato a spaglio o più frequentemente a file. Successivamente il seme viene ricoperto con della sostanza organica o sabbia (protezione dai predatori e/o per necessità di germinazione). Questa parte del vivaio si chiama semenzaio, ovvero dove si semina.\\
    Le piante rimangono in campo 1-4 anni, in funzione alla velocità di accrescimento (es. abete rosso 1-2, pino cembro 3-4).\\
    Dopo il periodo, le piante vengono svellite (rimosse dal semenzaio) e poste in una seconda aiuola (piantonaio): questo cambiamento serve per allargarle.\\
    In questo modo, le piantine a radice nuda passano solitamente un periodo in semenzaio ed un periodo in trapianto. Ogni pianta può avere nell'etichetta il codice S + T, seguiti dal numero di anni che sono stati in semenzaio e/o in trapianto. Ci sono altre due aree del vivaio: la bastardiera (dove finiscono le piantine invendute ed abbastanza grosse) e la zona a riposo (che ospitano le piante malate). 
    \item per le piante a pane di terra: normalmente non ha un campo, ma solitamente un pavimento in cemento sopraelevato dal suolo (in modo da lavorare ad altezza uomo). Quasi sempre, sono posti in serre, con l'impianto di irrigazione per aspersione.
\end{itemize}
In comune tra i vivai a pane di terra ed a radice nuda sono le infrastrutture. I vivai dovrebbero essere:
\begin{multicols}{2} 
\begin{itemize}[nosep]
    \item in zone piane, per la meccanizzazione;
    \item in aree consone alle specie che si andranno ad allevare; 
    \item facilmente accessibili dai mezzi di trasporto;
    \item recintati, per la protezione dagli animali selvatici;
    \item dotato di acqua, sia questa di fiume o di falda (pozzo).
\end{itemize}
\end{multicols}
L'acqua, data la bassa temperatura di acquisizione, potrebbe generare shock durante l'irrigazione: per questo motivo dev'essere in qualche modo riscaldata.\\
L'irrigazione può essere fatta per:
\begin{itemize}[nosep]
    \item aspersione: costo dell'impianto di irrigazione, permette di difendersi dalle gelate tardive;
    \item a goccia: sistema più efficiente, essendo puntuale per ogni pianta. Ha un alto costo di impianto;
    \item sommersione: avviene l'allagamento del terreno;
    \item infiltrazione laterale: mediante la creazione di canalette laterali. 
\end{itemize}
I vivai hanno diversi locali per il rimesso dei macchinari, uffici, luoghi per la conservazione del seme, etc.\\
Pistoia è l'area più vocata al vivaismo in Italia.\\
I vivai si sono ridotti di numero nel tempo, anche perché le attività di oggi prevedono la rinnovazione naturale.\\
La previsione del rimboschimento dev'essere comunicata al vivaio con sufficiente preavviso, in modo che il materiale forestale possa essere preparato.\\
Tendenzialmente, la radice nuda si utilizza per specie fittonanti ed a rapido accrescimento (c'è il rischio di spiralatura delle radici se messe in vaso).
\subsection{Pianificazione del rimboschimento}
\subsubsection{Sesti}
Possono essere geometrici: quadrato, rettangolare, quittonce e settonce. Sono quelli tradizionali del rimboschimento (in bosco i primi due sono quelli più usati).\\
Questi rimboschimenti sono di facile pianificazione/realizzazione, ma anche di facile riconoscimento (essendo regolari).\\
In questi impianti, dove le distanze sono regolari e le piantine sono simili, la crescita è uguale per tutte e quindi c'è molta instabilità.
\paragraph{In pianura}
Ultimamente, soprattutto in pianura, si tende a fare il rimboschimento a gruppi, che va bene anche per i sottoimpianti. Solitamente, la rinnovazione è densa, e perciò si pianta il numero finale di gruppi (es. ``se voglio 100 piante, pianto 100 gruppi''): ogni gruppo ha un certo numero di piantine (es. 8-12 cadauno).\\
Normalmente, le piantine vengono piantate a 30-50 cm l'una dall'altra; in questo modo si va a simulare l'iniziale densità della rinnovazione (facilitazione): alla fine rimarrà solamente una piantina.\\
La ripulitura è più veloce, poiché si deve solamente ripulire attorno al gruppo, senza entrarci.\\
Il numero di piantine è solitamente di 2000-2500 piantine ad ettaro. Questo sistema si presta bene per tutte le specie che crescono velocemente (es. le querce in pianura).
\paragraph{In montagna}
In montagna, dove la struttura del bosco è normalmente a collettivi, ed un rimboschimento geometrico normale è destinato a fallire (la facilitazione è molto più importante della competizione), vengono svolti così rimboschimenti a collettivi:
\begin{itemize}[nosep]
    \item viene individuata la microstazione favorevole (es. un microdosso, una ceppaia);
    \item intorno viene costruito il collettivo. Inizialmente venivano piantate le piante in modo regolare, a distanze troppo elevate. Oggi si creano dei microcollettivi pari al numero delle piante che occuperanno quel collettivo (es. se nel collettivo rimarranno 10 piante, si piantano 10 microcollettivi).\\
    Nella realtà: all'interno del microcollettivo la facilitazione è importante e la chiusura delle chiome è veloce. La distanza tra i microcollettivi è inferiore alla distanza media del raggio di chioma di una pianta adulta, in questo modo i diversi microcollettivi andranno a chiudersi subito; infine, si chiuderà anche il collettivo. La distanza tra due collettivi è sempre tale per cui le piante non si tocchino mai, anche alla loro massima grandezza.\\
    Anche in questo caso le piante usate per il rimboschimento sono 2000-2500 per ettaro. 
\end{itemize} 
\subsubsection{Operazione di rimboschimento}
Il periodo dell'anno per svolgere il rimboschimento dipende dal materiale:
\begin{itemize}[nosep]
    \item pianta a radice nuda: non può essere fatta durante la stagione vegetativa (eccessivo shock da trapianto). Può essere in autunno (dopo la fine della stagione vegetativa), in inverno o in primavera (prima dell'inizio della stagione vegetativa).\\
    Occorre evitare il pieno inverno (basse temperature e neve).\\
    L'autunno e la primavera sono buoni periodi per piantare. La primavera sarebbe la scelta tecnicamente migliore, perché subito dopo il trapianto si entra in periodo vegetativo; il problema è la durata delle operazioni di rimboschimento, che potrebbero portare ad un ritardo sulla fine delle operazioni. Per questo motivo, l'autunno è il miglior periodo dal punto di vista operativo, perché c'è un maggiore franco di lavoro (``e male che vada c'è la primavera"). 
    \item pianta in pane di terra: il ragionamento è il medesimo della radice nuda; in questo caso però è possibile anticipare e/o posticipare leggermente verso la stagione vegetativa.\\
    L'inverno e l'estate sono da evitare, anche perché le piante continuano comunque a crescere all'interno del pane di terra in attesa del trapianto (rischio cioè di piantare materiale spiralato). 
\end{itemize}